% !TEX root = ../main_zh.tex
\chapter{凯恩斯交叉与乘数}
\label{ch:keynes}

前几章的增长模型告诉我们经济长远会走向何处。在索洛模型与新古典框架中（第~\ref{ch:solow} 章与第~\ref{ch:neoclassical} 章），以及每当我们写下
\[
Y = F(K, L, \mathrm{TFP})
\]
时，产出都是由供给一侧决定的：取决于资本存量、劳动力数量，以及把二者转化为产品的技术水平。对于\emph{长期}而言，这是正确的图景——此时价格与工资已有足够时间调整，经济运行在其生产潜力之上。但它对另一些问题闭口不谈：产出为什么会逐年起落，衰退为什么发生，政府又为什么会认为一次减税就能把经济拉出低谷。要回答这些问题，我们需要一个\emph{短期}模型——在其中，决定产出的不是经济能生产多少，而是买家愿意花多少钱。

本章发展的正是这样一个最简单的模型，它出自凯恩斯之手。其组织性恒等式，是我们在国民收入核算中早已熟悉的产出需求侧分解（第~\ref{ch:accounts} 章）：
\[
Y = C + I + G + \mathrm{NX},
\]
只是如今要把它读作关于\emph{计划}支出的陈述，而非已实现的会计记录。我们将逐项说明消费、投资、政府购买与净出口这四个分量如何依赖于收入，把它们加总为一条计划支出曲线，再施加“计划支出等于产出”这一要求。由此得到的便是\emph{凯恩斯交叉}（Keynesian cross）——一张其交点钉住均衡产出的图。它最具分量的教益是\emph{乘数}（multiplier）：由于任何一笔额外支出都会变成某人的收入并被再次部分花出，对需求的一点小小推动会让产出移动得比这点推动本身更多。

关于这个模型究竟是什么，要先说一句提醒的话。凯恩斯的框架在本质上是\emph{静态}的，它没有清晰的“未来”概念：消费只对\emph{当期}收入作出反应，投资被处理成与当期产出无关，而利率、终生财富之类的跨期因素，要么被压下不谈，要么被统统折叠进一个行为参数里。这是一个实实在在的局限，后面的章节（第~\ref{ch:islm} 章的 IS--LM 模型与第~\ref{ch:phillips} 章的动态分析）会把这里缺失的东西补回去一部分。然而模型的粗糙本身也正是它历久弥新的源头：它抓住了一种力量——由当期收入产生的支出——而短期数据恰恰最强烈地呈现出这种力量，并且它把这种力量装进了一个足够简单、可以用来推理的形式之中。

\section{产出的两种时间视野}

值得把贯穿下文的那一组对照明白地说出来。长期之下，产出是供给侧的量：$Y = F(K, L, \mathrm{TFP})$ 由要素存量与技术状态固定，需求至多只能决定价格。短期之下，价格黏滞、要素可能闲置，因果方向便反了过来。需求的水平——家庭、企业、政府与外国想要购买多少——决定了企业实际生产多少，从而决定产出与就业。

会计恒等式 $Y = C + I + G + \mathrm{NX}$ 在\emph{每一}期都按定义成立，因为国民收入核算记录的是已实现的购买。凯恩斯的做法，是把右端重新诠释为\emph{计划}支出或\emph{意愿}支出，再追问：什么时候这些计划与它们所预设的产出彼此自洽？由于支出对应着需求，我们就把计划支出当作经济的总需求，
\[
\mathrm{AD} \;=\; C + I + G + \mathrm{NX},
\]
并不再多加辩白地把二者视为等同。这一等同并不绝对严谨——在更完整的模型里，总需求还会对价格水平和利率作出反应——但在短期中，需求的波动如此之大、又远比潜在产出来得迅速，以至于就本章的目的而言，我们大可把产出的短期变动当作需求的变动，无须再区分一条短期曲线与一条长期曲线。

\section{计划支出的各个分量}

现在我们为四个分量各写下一个行为方程，逐项搭建起计划支出曲线。全章中，$Y$ 记收入（也即产出），$T$ 记扣除转移支付后的净税收，因而 $Y - T$ 就是\emph{可支配收入}——家庭手中可供花费或储蓄的那部分收入。

\subsection{消费}

一个家庭消费多少由什么决定？原则上，决定因素很多：它的当期收入、它预期的未来收入（极端而言是它的整个\emph{终生}收入）、它积累的财富（一个存量，与消费正相关），以及利率、对未来的贴现率这类跨期项。凯恩斯的模型有意忽略了这其中的绝大部分。由于它没有真正的动态结构——没有可以享用所储蓄资源的未来期——它无法把利率或终生收入表示为独立的力量。消费被设定为只依赖于\emph{当期}可支配收入：

\defn{凯恩斯消费函数}{
家庭消费是当期可支配收入的线性函数，
\[
C \;=\; C_0 + \mathrm{MPC}\,\pr{Y - T},
\]
其中 $C_0 > 0$ 是\emph{自发型消费}（autonomous consumption）——即可支配收入为零时仍会发生的消费量；$\mathrm{MPC} \in (0,1)$ 是\emph{边际消费倾向}（marginal propensity to consume），即每多一元可支配收入中被花掉的比例。$\mathrm{MPC}\,(Y-T)$ 一项是\emph{引致消费}（induced consumption），即由收入驱动的那部分。
}

关键的建模约定是：\emph{所有变量都是当期的}——方程里没有预期未来收入这一项。未来、财富或利率在塑造行为时所起的任何作用，都被整体吸收进 $\mathrm{MPC}$ 这一个参数里。一个有耐心、富有或对未来乐观的家庭，无非是拥有一个不同的 $\mathrm{MPC}$；至于这个数字从何而来，模型并不交代。

\rmkb[一个粗糙却合乎数据的模型]{这条消费函数显然是不完整的——它略去了财富、预期和利率，而这些都是更丰富的理论会纳入的因素。凯恩斯本人并不否认这一点。他的辩护是经验性的：这条函数是从观察、而非从第一原理中搭起来的，而可支配收入事实上正是能解释总消费变动最多的那个单一因素。我们应当把 $C = C_0 + \mathrm{MPC}\,(Y-T)$ 读作\emph{基础}模型，再随应用所需向其追加更多项（一个利率效应、一个财富效应）。}

\rmkb[减税与中国的边际消费倾向]{$\mathrm{MPC}$ 的大小有着直接的政策含义，而且这些含义因制度而异。中国的税收减免历来主要通过削减\emph{企业}税来实现——税收的大头也正是由企业贡献的——而非通过对个人的转移支付；向家庭直接发放现金补贴的做法并不常见。再叠加上相对偏低的家庭 $\mathrm{MPC}$，这意味着一笔给定的减税往往只能微弱地传导到消费上。相比之下，以投资为基础的刺激则能在两个边际上同时发力：短期内它拉动就业与需求，长期内它又能通过增添生产能力而改变供给一侧。}

储蓄不过是可支配收入中未被消费的那一部分，因此消费函数立刻就蕴含了一条\emph{储蓄函数}：
\[
S \;=\; \pr{Y - T} - C \;=\; -C_0 + \pr{1 - \mathrm{MPC}}\pr{Y - T}.
\]
斜率 $1 - \mathrm{MPC}$ 是\emph{边际储蓄倾向}：每一元可支配收入非花即存，所以这两个倾向之和恒为一。

\subsection{投资}

一家企业只在某个投资项目有利可图时才会去做它。自然的判据是净现值：把项目未来的收益流与成本流贴现回当下，当前者超过后者时就投资，
\[
\mathrm{NPV}(I) \;=\; \sum_{t=1}^{\infty} \frac{\text{Income}_t}{\pr{1 + r}^{t}} \;-\; \sum_{t=1}^{\infty} \frac{\text{Cost}_t}{\pr{1 + r}^{t}}.
\]
麻烦出在收益一侧。未来收益高度不确定，且大多难以量化刻画；在实践中，这一决策靠的更接近“直觉”而非计算——也就是凯恩斯所说的著名的\emph{动物精神}（animal spirits），他以动物追逐猎物的本能冲动作类比。成本一侧则更易于把握。投资必须融资，而融资成本随利率 $r$ 上升，于是利率越高，能跨过门槛的项目就越少。因此我们把投资建模为利率的减函数，
\[
I = I(r), \qquad I'(r) < 0,
\]
而不去认定某个具体的函数形式——凯恩斯本人也没有提出任何具体形状。

\state{投资不依赖于当期收入}{
在凯恩斯模型中，投资仅是利率的函数：$I = I(r)$，关于 $r$ 递减。它\emph{不}依赖于当期产出、家庭收入或财富。在这个静态模型里，这是一个有意为之的简化，目的是让 $I$ 与 $Y$ 的决定因素彼此分开。
}

\rmkb[现实中的保留意见]{“投资独立于当期收入”这一假设是一个干净的理论立场，而非对现实世界的字面描述。现实当中，投资与产出强相关：它会对未来需求的预期、对政府政策作出反应，也与决策者的财富和信心、与内部资金的可得性相关。把 $I$ 当作外生于 $Y$，最好读作一种“一次只孤立一条渠道”的办法；一旦我们让投资对利率作出反应，而这个利率又由第~\ref{ch:islm} 章的 IS--LM 模型与产出联合决定，这一假设便被放松了。}

\subsection{政府购买}

政府对商品与服务的购买 $G$，被视为某种政治与行政过程的产物，而非对当期经济变量的反应。它由决策定下，不由模型定下，因此我们把 $G$ 当作\emph{外生}的：一个由政策制定者选定的常数。税收 $T$ 同样处理，当作一个由模型之外设定的政策工具。

\subsection{净出口}

净出口是出口减去进口，
\[
\mathrm{NX} = X - M.
\]
出口 $X$ 是外国人对本国产品的购买，因而依赖于\emph{外国}收入，并与之正相关；从本国模型的视角看，它是外生的。进口 $M$ 是本国对外国产品的购买，随本国可支配收入上升，
\[
M = \mathrm{MPM}\,\pr{Y - T},
\]
其中 $\mathrm{MPM}$ 是\emph{边际进口倾向}（marginal propensity to import），即每多一元可支配收入中花在外国产品上的比例。在大部分分析中，把这种收入依赖性压下、同样把 $\mathrm{NX}$ 当作一个外生常数，往往更方便，也更常见；我们在推导乘数时就会这样做，并在一则注记中说明对收入敏感的情形会如何修正结果。

\section{计划支出与凯恩斯交叉}

把四个分量加起来，便得到作为收入函数的计划支出曲线。把依赖收入的各项明确写出来，
\[
\mathrm{AD} \;=\; C + I + G + \mathrm{NX}
\;=\; \underbrace{C_0 + I(r) + G + X}_{\text{自发部分}}
\;+\; \underbrace{\pr{\mathrm{MPC} - \mathrm{MPM}}}_{\text{斜率}}\pr{Y - T}.
\]
这条曲线随收入递增，但斜率严格小于一：收入每多一元，计划支出只上升 $\mathrm{MPC} - \mathrm{MPM} < 1$，因为这多出的一元里有一部分被储蓄、一部分外漏到了国外。在一张纵轴为计划支出、横轴为收入 $Y$ 的图中，这是一条向上倾斜、但比 $45^\circ$ 线更平缓的直线。（若把进口折进外生常数里，斜率就只是 $\mathrm{MPC} < 1$；真正要紧的，只是它为正且小于一。）

均衡要求计划彼此自洽：企业生产的恰好是买家想买的，从而计划支出等于产出，不存在非意愿的存货堆积或消耗。

\defn{短期均衡}{
当计划总支出等于产出时，经济处于短期均衡，
\[
\mathrm{AD} \;=\; Y,
\]
也即计划支出曲线与“支出等于收入”的那条 $45^\circ$ 线相交之处。
}

$45^\circ$ 线是所有“需求量等于产出量”的点的轨迹；它正是界定均衡的约束条件。由于计划支出曲线的斜率小于一，而 $45^\circ$ 线的斜率恰为一，两条线必定恰好相交一次。它们的交点就是均衡产出水平。这张图便是\emph{凯恩斯交叉}，如图~\ref{fig:macro-12} 所示。

\begin{figure}[h]
\centering
\begin{tikzpicture}[scale=1.0]
  % axes
  \draw[-Latex] (0,0) -- (8.6,0) node[right] {$Y$};
  \draw[-Latex] (0,0) -- (0,7.4) node[above] {$\mathrm{AD}$};
  % 45-degree line (equilibrium condition AD=Y, slope 1, steeper) through origin
  \draw[line width=1.2pt] (0,0) -- (7.1,7.1);
  \node[above left=1pt and -2pt] at (7.1,7.1) {$45^{\circ}$};
  % planned aggregate expenditure line: positive intercept, slope < 1 (flatter)
  \draw[semithick] (0,1.6) -- (8.0,6.0);
  \node[below right=-1pt and 1pt] at (8.0,6.0) {$\mathrm{AD} = C+I+G+\mathrm{NX}$};
  % equilibrium: intersection of y=x with y=1.6+0.55x
  % slope of AD line = (6.0-1.6)/8.0 = 0.55 ; eq: x = 1.6/(1-0.55) = 3.5556
  \pgfmathsetmacro{\xeq}{1.6/(1-0.55)}
  \fill (\xeq,\xeq) circle (1.7pt);
  \draw[densely dashed] (\xeq,\xeq) -- (\xeq,0) node[below] {$Y_{0}$};
\end{tikzpicture}

\caption{凯恩斯交叉：均衡产出出现在计划支出曲线（斜率小于一）与“支出等于收入”的 $45^\circ$ 线相交之处。}
\label{fig:macro-12}
\end{figure}

\subsection{均衡的稳定性}

这个交点不仅仅是方程恰好相抵的一点；它还是一个\emph{稳定}的休止点，经济从两侧都会被推向它。其调整机制经由存货与企业的生产决策运作。设产出一时高于均衡水平，$Y_1 > Y_0$。在 $Y_1$ 处，计划支出线位于 $45^\circ$ 线\emph{之下}，于是计划支出不及产出：买家想要的比企业已生产的少。商品作为未售存货堆积起来，企业便缩减生产，产出向 $Y_0$ 漂回。反过来，若产出过低，$Y_2 < Y_0$，计划支出超过产出，存货被消耗，企业扩大生产，产出又升回 $Y_0$。从两侧看，经济都收敛到那个交点，如图~\ref{fig:macro-13} 所示。

\begin{figure}[h]
\centering
\begin{tikzpicture}[scale=1.0]
  % --- geometry -------------------------------------------------------
  % 45-degree line: AD = Y  (through the origin, slope 1)
  % AD line:        AD = a + b*Y, intercept a>0, slope b<1 (flatter)
  % cross at Y0 where a + b*Y0 = Y0  =>  Y0 = a/(1-b)
  \def\a{1.3}\def\b{0.45}
  \pgfmathsetmacro{\Yz}{\a/(1-\b)}        % Y0
  \pgfmathsetmacro{\Yone}{3.4}            % Y1 > Y0 (pulled nearer Y0, still right of it)
  \pgfmathsetmacro{\Ytwo}{1.35}           % Y2 < Y0
  \pgfmathsetmacro{\xmax}{8.2}
  \pgfmathsetmacro{\ymax}{6.6}
  % helper: AE(Y) and 45(Y)
  % --- axes -----------------------------------------------------------
  \draw[-Latex] (0,0) -- (\xmax,0) node[right] {$Y$};
  \draw[-Latex] (0,0) -- (0,\ymax) node[above] {$\mathrm{AD}$};
  % --- the two lines --------------------------------------------------
  % 45-degree line (steeper, primary): thicker
  \pgfmathsetmacro{\fortyfive}{\ymax-0.3}
  \draw[line width=1.2pt] (0,0) -- (\fortyfive,\fortyfive)
        node[above left=-2pt and -1pt] {$45^{\circ}$};
  % AD line (flatter): thinner
  \pgfmathsetmacro{\aeEnd}{\a+\b*\xmax}
  \draw[semithick] (0,\a) -- (\xmax,\aeEnd);
  \node[right] at (\xmax,\aeEnd) {$\mathrm{AD}=C+I+G+\mathrm{NX}$};
  % --- equilibrium point at Y0 ---------------------------------------
  \fill (\Yz,\Yz) circle (1.6pt);
  % --- dashed guides from each output level up to its line -----------
  % at Y0: up to the cross
  \draw[densely dashed] (\Yz,0) -- (\Yz,\Yz);
  % at Y1 (right of Y0): the 45-degree line is higher -> reach the 45 line
  \draw[densely dashed] (\Yone,0) -- (\Yone,\Yone);
  % at Y2 (left of Y0): the AD line is higher -> reach the AD line
  \pgfmathsetmacro{\aeTwo}{\a+\b*\Ytwo}
  \draw[densely dashed] (\Ytwo,0) -- (\Ytwo,\aeTwo);
  % --- output-level labels on the axis -------------------------------
  \node[below] at (\Ytwo,-0.05) {$Y_2$};
  \node[below] at (\Yz,-0.05)   {$Y_0$};
  \node[below] at (\Yone,-0.05) {$Y_1$};
  % --- adjustment arrows toward Y0 (drawn just below the axis) -------
  % at Y1: AD<Y, firms cut output -> leftward arrow toward Y0
  \draw[-Latex,thick] (\Yone,-0.85) -- (\Yz+0.35,-0.85);
  \node[below=1pt] at ({(\Yone+\Yz)/2+0.2},-0.95) {\footnotesize $\mathrm{AD}<Y$};
  % at Y2: AD>Y, firms expand output -> rightward arrow toward Y0
  \draw[-Latex,thick] (\Ytwo,-0.85) -- (\Yz-0.35,-0.85);
  \node[below=1pt] at ({(\Ytwo+\Yz)/2-0.1},-0.95) {\footnotesize $\mathrm{AD}>Y$};
\end{tikzpicture}

\caption{凯恩斯交叉均衡的稳定性：在 $Y_0$ 之上，计划支出不及产出，企业缩减；在 $Y_0$ 之下，支出超过产出，企业扩张；产出从两侧都收敛到 $Y_0$。}
\label{fig:macro-13}
\end{figure}

\rmkb[均衡产出未必等于潜在产出]{长期来看，凯恩斯交叉的均衡本应与潜在产出重合——也就是供给所决定的那个水平 $Y = F(K,L,\mathrm{TFP})$。短期之下却没有这样的保证。经济可以稳定地停在潜在产出之下，伴随闲置产能与失业劳动；也可以停在其上，经济过热。模型给出的均衡与长期模型给出的潜在产出之间的这道缺口，正是短期需求管理之所以有理由存在的依据。}

\section{乘数}

两条线的斜率相对关系所确定的，不止是它们会相交这一点。它还支配着当计划支出曲线移动时，交点会\emph{移动多远}。设某项自发支出上升，使整条曲线向上平移了 $\Delta D$。由于这条曲线比 $45^\circ$ 线更平缓，交点在横向上走过的距离 $\Delta Y$ 会超过引起它的纵向平移 $\Delta D$：自发支出的小幅增加，换来了产出更大幅度的增加。这种放大就是\emph{乘数效应}，而平移被放大的倍数就是\emph{乘数}。其几何关系如图~\ref{fig:macro-14} 所示。

\begin{figure}[h]
\centering
\begin{tikzpicture}[scale=1.0,>=Latex]
  % --- geometry parameters ---
  % AE line: AD = a + m*Y, with m<1 (flatter than 45-degree line)
  \def\m{0.5}            % MPC slope of the AE line
  \def\al{1.2}           % intercept of the LOWER (initial) AE line
  \def\dD{1.2}           % vertical upward shift Delta D
  % upper intercept = al + dD
  % equilibria: Y = a/(1-m)
  \pgfmathsetmacro{\Yz}{\al/(1-\m)}          % Y_0
  \pgfmathsetmacro{\Yo}{(\al+\dD)/(1-\m)}    % Y_1
  \pgfmathsetmacro{\ADz}{\Yz}                 % AD at Y_0 (on 45 line)
  \pgfmathsetmacro{\ADo}{\Yo}                 % AD at Y_1 (on 45 line)
  \def\xmax{8.4}
  \def\ymax{7.2}

  % --- axes ---
  \draw[-Latex] (0,0) -- (\xmax,0) node[right] {$Y$};
  \draw[-Latex] (0,0) -- (0,\ymax) node[left] {$\mathrm{AD}$};

  % --- 45-degree line ---
  \def\Lf{6.6}
  \draw[line width=1.1pt] (0,0) -- (\Lf,\Lf);
  \node[above left=-1pt and 0pt] at (\Lf,\Lf) {$45^{\circ}$};

  % --- AE lines (planned expenditure), thicker = primary feasible lines ---
  \def\xR{7.6}
  % lower (initial) AE line
  \pgfmathsetmacro{\yLR}{\al+\m*\xR}
  \draw[line width=1.3pt] (0,\al) -- (\xR,\yLR);
  % upper (shifted) AE line
  \pgfmathsetmacro{\yUR}{\al+\dD+\m*\xR}
  \draw[line width=1.3pt] (0,\al+\dD) -- (\xR,\yUR);
  \node[right] at (\xR,\yUR) {$C+I+G+\mathrm{NX}$};

  % arrow showing the upward shift of the AE line
  \draw[-Latex] (\xR-1.0,{\al+\m*(\xR-1.0)}) -- (\xR-1.0,{\al+\dD+\m*(\xR-1.0)});

  % --- equilibria ---
  \fill (\Yz,\ADz) circle (1.6pt);
  \fill (\Yo,\ADo) circle (1.6pt);

  % --- dropped guide lines to the x-axis ---
  \draw[densely dotted] (\Yz,\ADz) -- (\Yz,0) node[below] {$Y_0$};
  \draw[densely dotted] (\Yo,\ADo) -- (\Yo,0) node[below] {$Y_1$};

  % --- Delta D : vertical gap between the two AE lines (left side) ---
  \def\xD{1.4}
  \pgfmathsetmacro{\yDlo}{\al+\m*\xD}
  \pgfmathsetmacro{\yDhi}{\al+\dD+\m*\xD}
  \draw[<->,green!55!black,line width=0.9pt] (\xD,\yDlo) -- (\xD,\yDhi);
  \node[left=1pt,orange!85!black] at (\xD,{(\yDlo+\yDhi)/2}) {$\Delta D$};

  % --- Delta Y : horizontal gap between Y_0 and Y_1 on the axis ---
  \def\yDY{1.0}
  \draw[<->,green!55!black,densely dotted,line width=1.0pt] (\Yz,\yDY) -- (\Yo,\yDY);
  \node[above,orange!85!black] at ({(\Yz+\Yo)/2},\yDY) {$\Delta Y$};
\end{tikzpicture}

\caption{乘数：计划支出曲线向上平移 $\Delta D$ 时，均衡产出上升 $\Delta Y > \Delta D$，原因是这条曲线比 $45^\circ$ 线更平缓。}
\label{fig:macro-14}
\end{figure}

要在代数上求出乘数，就把 $\mathrm{NX}$ 当作外生常数，并对计划支出曲线
\[
\mathrm{AD} \;=\; C_0 + \mathrm{MPC}\,\pr{Y - T} + I + G + \mathrm{NX}
\]
施加均衡条件 $\mathrm{AD} = Y$。令 $\mathrm{AD} = Y$ 并解出 $Y$：
\[
\ba
Y &= C_0 + \mathrm{MPC}\,\pr{Y - T} + I + G + \mathrm{NX}\\
Y - \mathrm{MPC}\,Y &= C_0 - \mathrm{MPC}\,T + I + G + \mathrm{NX}\\
\pr{1 - \mathrm{MPC}}\,Y &= C_0 - \mathrm{MPC}\,T + I + G + \mathrm{NX},
\ea
\]
即得均衡产出的简化式解，
\begin{equation}
Y \;=\; \frac{1}{1 - \mathrm{MPC}}\,\pr{C_0 - \mathrm{MPC}\,T + I + G + \mathrm{NX}}.
\label{eq:keynes-Y}
\end{equation}
现在，每个乘数都可以读作 $Y$ 对相应外生变量的偏导数。

\thm{凯恩斯乘数}{
在均衡式~\eqref{eq:keynes-Y} 中，产出对各外生分量的反应为：
\[
\ba
\text{支出乘数：} \quad & \pd{Y}{G} = \pd{Y}{I} = \pd{Y}{C_0} = \frac{1}{1 - \mathrm{MPC}} \;>\; 1,\\[4pt]
\text{税收乘数：} \quad & \pd{Y}{T} = \frac{-\,\mathrm{MPC}}{1 - \mathrm{MPC}} \;<\; 0,\\[4pt]
\text{平衡预算乘数：} \quad & \dd{Y}{G}\bigg|_{\d T = \d G} = 1.
\ea
\]
对任意 $\mathrm{MPC} \in (0,1)$，支出乘数都大于一；税收乘数为负，其\emph{绝对值} $\dfrac{\mathrm{MPC}}{1-\mathrm{MPC}}$ 大于一当且仅当 $\mathrm{MPC} > \tfrac12$。
}

\pf{
支出乘数由对~\eqref{eq:keynes-Y} 求导得到：由于 $C_0$、$I$、$G$、$\mathrm{NX}$ 都以系数 $1/(1-\mathrm{MPC})$ 进入，它们各自的乘数都是 $1/(1-\mathrm{MPC})$，而这个值因 $0 < \mathrm{MPC} < 1$ 而大于一。其经济内涵就是一轮接一轮的“漏出”故事：多出的一元自发支出变成某人的一元收入，其中比例 $\mathrm{MPC}$ 被再次花出，再花出的部分又有 $\mathrm{MPC}$ 被花出，如此往复。引致出来的产出是一个等比级数
\[
1 + \mathrm{MPC} + \mathrm{MPC}^2 + \cdots = \frac{1}{1 - \mathrm{MPC}}.
\]
对~\eqref{eq:keynes-Y} 关于 $T$ 求导，得税收乘数 $-\mathrm{MPC}/(1-\mathrm{MPC})$，它为负，因为增税削减可支配收入、从而削减支出。它的绝对值恰好比支出乘数小一个因子 $\mathrm{MPC}$：税收的变化要先穿过消费函数才能作用于需求，所以它的第一轮效应是 $\mathrm{MPC}\,\Delta T$，而不是整个 $\Delta G$。解 $\mathrm{MPC}/(1-\mathrm{MPC}) > 1$ 得 $\mathrm{MPC} > 1 - \mathrm{MPC}$，即 $\mathrm{MPC} > \tfrac12$。

至于平衡预算乘数，令支出与税收同时增加同一数额，$\Delta G = \Delta T \equiv \Delta$。由刚求得的两个乘数，
\[
\Delta Y = \frac{1}{1-\mathrm{MPC}}\,\Delta G - \frac{\mathrm{MPC}}{1-\mathrm{MPC}}\,\Delta T
= \frac{1 - \mathrm{MPC}}{1 - \mathrm{MPC}}\,\Delta = \Delta.
\]
一笔完全由增税来融资的政府购买增加，使产出恰好上升等于这笔增加的数额：平衡预算乘数为一。
}

支出乘数与税收乘数之差背后的直觉值得多花些笔墨，因为它支撑着整场财政政策辩论。$G$ 或 $I$ 的增加会把计划支出曲线整额地向上推，因而被整个因子 $1/(1-\mathrm{MPC})$ 放大。减税则相反，它只能间接地抬高支出：家庭拿到了多出的可支配收入，却只花掉其中的比例 $\mathrm{MPC}$，把其余存了起来。于是曲线只向上平移 $\mathrm{MPC}\,\abs{\Delta T}$，所以就同样一元钱而言，税收是更弱的工具。平衡预算的结果随之道出一件惊人的事：哪怕是一笔完全靠增税来支付的刺激，也并非中性——它仍然会一比一地抬高产出，因为支出注入足额起作用，而税收的抽离却被较低的储蓄部分吸收掉了。

\rmkb[对收入敏感的净出口]{若保留进口对收入的敏感性、不把它折进常数，即 $\mathrm{NX} = X - \mathrm{MPM}(Y-T)$，则均衡条件变为 $Y = C_0 + \mathrm{MPC}(Y-T) + I + G + X - \mathrm{MPM}(Y-T)$，解之得
\[
Y = \frac{1}{1 - \mathrm{MPC} + \mathrm{MPM}}\,\pr{C_0 - (\mathrm{MPC}-\mathrm{MPM})T + I + G + X}.
\]
现在支出乘数是 $1/(1 - \mathrm{MPC} + \mathrm{MPM})$，比之前更小：进口是一笔额外的漏出，所以每一轮再支出失去的不只是被储蓄的那部分，还有花到国外的那部分。定性的教益——一个大于一的正支出乘数、一个绝对值更小的负税收乘数——则保持不变。}

\state{乘数为何对政策重要}{
乘数告诉我们：支出、税收或投资变动所产生的效果，并不局限于其自身的规模——自发需求的一点小变化撬动了产出的一个更大变化。在评估其中任何一项时，都必须算上这种被乘数放大的效果，而不只是直接效果。两个乘数之间的不对称还带有一层政策意味——在同样的预算成本下，政府购买的增加比同等规模的减税更能拉动产出。这正是为什么当经济面临下行压力时，典型的应对是一种偏重支出的财政扩张。我们将在第~\ref{ch:fiscal} 章详谈财政政策。
}

然而凯恩斯交叉只讲了短期故事的一半。在计算投资时，我们一直把利率 $r$ 当作固定的，连带把整个自发部分也当作固定的。但利率本身是一个均衡价格，由货币市场决定。让 $r$ 变动，就会描出一整族均衡产出水平——每一个利率对应一个凯恩斯交叉均衡——而把这一族曲线与货币市场结合起来，正是第~\ref{ch:islm} 章 IS--LM 模型的任务。
